\chapter*{Introduction}
\addcontentsline{toc}{chapter}{Introduction}

Software designs keep changing. We replace the representation of data, factor
processing out into shared components, and split one system into several
services. In every case the question is what is preserved across the change.
A change of data representation may leave the displayed shipping address
unchanged while altering what happens when that address is updated. To compare
the two representations, we must specify both what is read and which operations
must be preserved. This distinction will recur throughout the paper.

AI-generated changes make these questions increasingly important. As more
candidate changes become available, we need grounds for deciding which ones
preserve the required structure, where local consistency fails to extend to
the whole, and how much freedom remains among the permitted changes. This
paper develops mathematical objects and theorems for making these decisions.
We call the resulting foundations Algebraic Architecture Theory (AAT).

The starting point is an explicit choice of what counts as structure and what
a change must preserve. We call this choice a \emph{reading}. For the order
example, a reading specifies whether the comparison concerns only the displayed
address or also the update and the payment data it must retain. AAT defines
architecture objects and their structure-preserving maps relative to such
choices. The following five questions explain what we want to do with these
objects before we describe how they are constructed.

\section*{Five questions}

\textbf{When do local states fit together?}
Suppose that three services record the time of the same event using their own
clocks. For each connection, we specify a fixed offset for converting a time
from one clock to the other. If conversion from the first clock through the
second and third and back to the first does not recover the original time,
no assignment of offsets to the three clocks can satisfy all these conversion
rules. Each service's records may be internally consistent, yet following the
connections around the cycle reveals a contradiction. Can we correct each
service's timestamps to obtain records satisfying every specified conversion
rule? The question is whether a global state exists that simultaneously
satisfies the relations required on overlaps (Chapter~\ref{chap:2}).

\textbf{When does a change of reading preserve a diagnosis?}
Suppose that we extract a list of shipping addresses from order data, omitting
internal payment data. Conditions involving only addresses can still be
checked from this list. But checking conditions on individual entries and
detecting discrepancies between services are different tasks. In the clock
example, inspecting each service separately does not reveal the discrepancy
around the cycle. If we group several services into a single region for
inspection, we must also retain the connection and conversion data needed
for that diagnosis. Which information can we omit, and how can we repartition
the regions we inspect, while still detecting the same discrepancies that
corrections cannot remove (Chapter~\ref{chap:3})?

\textbf{When do different routes of change agree?}
When migrating an order-processing API, we want a direct migration from the
old version to the new one and a migration through an intermediate version to
yield the same order and the same behavior when it is read or updated.
Being able to undo each conversion does not by itself ensure agreement
between the two routes.
The order of operations matters too. We may split order processing and
inventory management into services and then extract the operations involved
in order cancellation, or extract those operations first and apply the same
splitting policy. Do the resulting designs agree on completion of the
cancellation and release of the inventory reservation, preserving the condition
that a completed cancellation has no outstanding reservation? What is needed
to put the data and operations obtained by the two routes in correspondence
and ensure that they satisfy the same conditions
(Chapters~\ref{chap:4}--\ref{chap:5})?

\textbf{Which changes preserve a comparison, and how can we distinguish them?}
Suppose that we extract the shipping-address update into a shared component.
The displayed address may remain the same even though an update mishandles
the payment data. We want moving old data to the new format and then updating
the address to yield the same order state, including payment data, as updating
first and then moving. There may be more than one invertible change of the internal representations
that preserves this correspondence between old and new.
Can the display and the pattern of calls alone tell us whether a candidate
change is compatible, or must we also inspect the internal values passed
through operations? We seek to identify the information needed, construct all
compatible changes, and describe the freedom in choosing them
(Chapters~\ref{chap:6}--\ref{chap:7}).

\textbf{When do local descriptions determine the whole?}
Suppose that the modernization of an intrabank transfer system is divided
among the account, payment, and accounting teams. Each team describes the
types of data it handles, its operations and required conditions, and the
correspondence between the old and new designs. For a given transfer, the
teams' descriptions must refer to the same transaction and amount, with
compatible completion conditions, if they are to combine into a single
transfer process. Their old-to-new correspondences must agree too: a
transaction migrated by the payment team and its record migrated by the
accounting team must still refer to the same transaction. What must be
described, and what must agree across teams, to assemble an overall design
and a single migration preserving its operations and conditions? When can
we also rule out distinct migrations that differ only where the descriptions
say nothing (Chapter~\ref{chap:8})?

These questions share several ingredients: objects, operations, local readouts
and their overlaps, and equations to be satisfied. We now describe the objects
that bring these ingredients together.

\section*{Architecture relative to a reading}

What facts are to be retained? An \emph{Atom} is a typed primitive fact, such
as the existence of a component or a named operation. A finite family of
Atoms together with relations is a \emph{configuration}. Adding structure
data, such as state sets and interpretations of operations, gives an
\emph{architecture object}. The equations that objects must satisfy are
specified as \emph{Laws}. These specify, for example, the required relationship
between reading and updating a state.

Which objects can arise under the permitted operations? From the extracted
Atom family we form a base object and generate a \emph{core}: it carries
the smallest family of objects closed under those operations. The core keeps
the choices of object formation, operations, and equations together. This
lets us ask what a change must preserve beyond the extracted facts alone.

To return to the three-service example, we also need to read information on
parts and compare it on overlaps. We choose local contexts, covers that
specify how the parts cover a whole, and coefficients that measure differences.
A core equipped with the covers, overlaps, coefficients, and local data is a
\emph{full geometry}. These are the objects on which the later chapters carry
out comparison and reconstruction.

The choices of vocabulary, extraction rules, object formation, equations,
operations, and locality together constitute a reading. Making the reading
explicit fixes the meaning of each preservation question. In the order
example, the choice to retain updates as well as displayed values changes
which maps are admitted. Chapter~\ref{chap:1} defines these choices and the
maps that compare them.

The \emph{Rising Sea} of the title comes from a metaphor with which
Grothendieck described his way of doing mathematics. A problem is treated not
as a hard shell to be cracked with a chisel, but as something that opens by
itself when immersed in a sea whose level rises quietly
(\cite[\href{https://webusers.imj-prg.fr/~leila.schneps/grothendieckcircle/RetSbis.pdf}{Part~III, note~122, ``La mer qui monte\ldots''}]{ReS};
the English phrase \emph{rising sea} follows the translation in
\cite{McLarty03}). Here the common objects allow different software questions
to be addressed by the same mathematical constructions. To explain how this
works, we next describe the connection to two independently defined semantics.

\section*{Correspondence with the semantics of CS}

In model synchronization, a \emph{lens} consists of states, views, reads, and
updates satisfying specified laws. In a protocol, named operations connect
states at control points, and observations record what can be read from those
states. We define these semantics independently in \S\S\ref{sec:1.8}--\ref{sec:1.9}.
For lenses, we construct Atoms and operations together with Laws that hold
if and only if the original laws hold (Proposition~\ref{prop:1.41}). Typed
morphisms recover the original semantics-preserving morphisms
(Proposition~\ref{prop:1.43}). For protocols, finite operation tables satisfying the declared path relations
and observation conditions determine realizations, and vertex maps satisfying
preservation conditions recover
semantics-preserving morphisms (Propositions~\ref{prop:1.37}
and~\ref{prop:1.38}).

These correspondences give a route from a software question to an AAT theorem
and back: construct the AAT input, prove the correspondence, apply the common
theorem, and interpret its conclusion in the original semantics. Each
correspondence states whether it gives an equivalence, an implication, or
preservation in one direction.

For example, consider the product lens used for the shipping address and
payment data. Of the four candidate changes that preserve the displayed
address and the unset payment value, only two also preserve updates. In the
worker protocol example, sixteen candidates preserve the correspondence of
control points, but only four preserve the handoff of contexts
(\S\ref{sec:7.8}). In both examples, an operation keeps an internal state
unchanged. Preserving that operation forces the changes of internal state at
its two ends to agree. The common classification theorem expresses this as
one permutation per connected component of the operation graph
(Theorem~\ref{thm:7.24}). Thus the same constraint explains both reductions in
the number of choices.

The correspondence also identifies, by a group isomorphism, the changes that
preserve a comparison in the original semantics and in the typed construction
(Theorem~\ref{thm:4.41}). Restricting the views of a lens also restricts its reads, updates, and lens
isomorphisms to the corresponding states (Proposition~\ref{prop:5.34}).
For protocols, realizations and semantics-preserving adapters restrict to
selected vertices, operations, and relations (Proposition~\ref{prop:5.36}). Section~\ref{sec:6.8}
examines which operations remain after normalization. At the end of the paper,
the same reconstruction principle yields unique extensions of morphisms from
compatible finite tables for lenses and protocols (Propositions~\ref{prop:8.26}
and~\ref{prop:8.28}).

\section*{Main constructions and their consequences}

The four parts organize the results as follows. The table gives a reading
guide; the paragraphs below state the principal conditions and conclusions.

\begin{xltabular}{\linewidth}{@{}>{\hsize=.7\hsize}L>{\hsize=1.15\hsize}L>{\hsize=1.15\hsize}L@{}}
\toprule
Part and chapters & Question & Results developed \\
\midrule
\endhead
I: Geometry and Local Consistency (Chapters~\ref{chap:1}--\ref{chap:2}) & What structure is compared, and when do local states fit together? & Construction of objects and geometry; obstructions to gluing and their relation to repair \\
\addlinespace[3pt]
II: Diagnosis and Transport (Chapters~\ref{chap:3}--\ref{chap:4}) & When does a diagnosis survive a change of reading, and when can structure be carried along a change? & Diagnostic invariance; universal transport and coherence of comparisons \\
\addlinespace[3pt]
III: Comparison and Normalization (Chapters~\ref{chap:5}--\ref{chap:6}) & When do two construction routes agree, and what does normalization retain? & Base-change comparisons; factorization and classification of invertibility \\
\addlinespace[3pt]
IV: Classification and Reconstruction (Chapters~\ref{chap:7}--\ref{chap:8}) & Which changes preserve a comparison, and when do local descriptions determine the whole? & Classification of compatible changes; recovery of objects and morphisms \\
\bottomrule
\end{xltabular}

\textbf{Geometry and local consistency (Part~I).}
Chapter~\ref{chap:1} constructs the core generated by a finite Atom family
and a reading. Its objects are exactly those reachable by finite sequences
of the permitted operations (Theorem~\ref{thm:1.18}). Contexts and covers
define a site; compatible restriction maps give a presheaf of rings and its
sheafification (\S\S\ref{sec:1.5}--\ref{sec:1.6}). The categories of
extraction, cores, and geometries are connected by projections that retain
the preceding levels of structure (Theorem~\ref{thm:1.31}).

Chapter~\ref{chap:2} expresses the satisfaction of Laws geometrically.
Under an algebraic presentation of the evaluation data and a localization
condition, simultaneous vanishing of the residuals is equivalent to
factorization through the space defined by the ideal of the equations
(Theorem~\ref{thm:2.9}). Differences of local states on overlaps determine a
degree-one \v{C}ech obstruction class. When local corrections give torsor
structures compatible with restriction and the states form a sheaf, vanishing
of this class is equivalent to the existence of a global state
(Theorem~\ref{thm:2.22}). Under the correspondence and completeness conditions
relating repair semantics to equations, a coefficient isomorphism identifies
their obstruction classes and connects vanishing with global repair
(Theorem~\ref{thm:2.49}, Corollary~\ref{cor:2.50}). An explicit
integer-coefficient example supplies an obstruction for the diagnostic
comparisons of the next chapter (\S\ref{sec:2.8}).

\textbf{Diagnosis and transport (Part~II).}
Chapter~\ref{chap:3} asks whether a diagnosis is preserved when the reading
changes. It constructs the coarsest resolution retaining all chosen Law
values (Theorem~\ref{thm:3.4}) and comparison maps between diagnostic
complexes. Under Condition~C on covers and coefficients, the comparison
induces an isomorphism on first cohomology (Theorem~\ref{thm:3.17}). Uniform
invariance, over all finite Law families for which both resolutions are
adequate, reduces to vanishing of kernels and cokernels indexed by nonempty
subsets of values. It is decidable when the required finite enumerations,
value tables, supports, and incidence data are computable
(Theorem~\ref{thm:3.22}, Proposition~\ref{prop:3.23}). The chapter also maps
the integer obstruction to the diagnosis and, under conditions on relation
components and common representatives, reflects vanishing
(Theorems~\ref{thm:3.37} and~\ref{thm:3.40}).

Chapter~\ref{chap:4} addresses a different question: carrying the structure
itself along a change. An exact change gives a bijective reindexing of
extracted Atoms. The transported core is characterized by unique
factorization of further changes (Theorem~\ref{thm:4.5}); transport to
geometry and its coherence with composition and projections are then
established (Theorems~\ref{thm:4.11}, \ref{thm:4.15}, and~\ref{thm:4.16}).
Specified comparisons may still differ from the canonical ones. For finite
comparison diagrams, the obstruction to removing these discrepancies
vanishes exactly when one can reselect edges to make all faces commute
(Theorem~\ref{thm:4.23}). Diagnosis and transport thus answer distinct
preservation questions about the structures constructed in Part~I.

\textbf{Comparison and normalization (Part~III).}
Chapter~\ref{chap:5} compares transporting structure with selecting an input
scope by pullback. It constructs the corresponding original input as a fiber
product and pulls back objects, operations, and Laws. For exact pointed
pullback squares, transport and pullback are related by an invertible
Beck--Chevalley comparison (Theorem~\ref{thm:5.11}). This is the precise
sense in which the two orders agree. For refinements that enlarge the
extracted content, backward transport instead requires reflection of
extraction on the realized locus (Theorem~\ref{thm:5.25}).

The chapter then builds two routes of full geometries from the same square,
finite comparison diagram, and geometry. Their generated comparison
$\beta=E\alpha$ combines an invertible comparison $\alpha$ with a
normalization $E$ chosen according to the diagnosis (\S\ref{sec:5.10}).
Chapter~\ref{chap:6} proves that, under the preservation conditions on
operations, residuals, invariants, and signatures, normalization defines an
idempotent morphism of full geometries (Theorem~\ref{thm:6.14}). The generated
comparison is invertible exactly when its normalization factor is the
identity (Theorem~\ref{thm:6.16}). Otherwise, the same comparison still gives
an isomorphism between the retained images in the idempotent completion
(Theorem~\ref{thm:6.12}). This distinguishes agreement after normalization
from recovery of the original structure.

\textbf{Classification and reconstruction (Part~IV).}
Chapter~\ref{chap:7} studies changes of the two endpoints that preserve a
chosen comparison. It separates two tasks: deciding whether a given change
is compatible, and constructing a compatible change from one specified after
normalization. Observations suffice for the first task exactly when the
kernel of the observation homomorphism is contained in the comparison-preserving
group (Theorem~\ref{thm:7.9}). For an invertible comparison and canonical
normalizations with a common presentation of the accompanying data, the
restriction homomorphism has a section. The compatible lifts of each fixed
normalized change form a torsor under its kernel
(Theorems~\ref{thm:7.19} and~\ref{thm:7.21}). Such lifts can exist even when
the normalized information does not determine compatibility of a given
change (\S\S\ref{sec:7.4}--\ref{sec:7.5}). For systems whose operations keep
hidden states, the classification by connected components also gives a split
short exact sequence onto the visible group (Theorem~\ref{thm:7.25}).

Chapter~\ref{chap:8} asks how to recover the objects and maps themselves
from local descriptions. The reconstruction principle has three requirements:
local readouts distinguish morphisms, compatible local morphisms can be
assembled, and local objects can be assembled (Theorem~\ref{thm:8.12}). The
main theorem proves these requirements for the primitive data of full
geometries, including types, operations, Laws, local geometry, and realization.
It gives an equivalence between the category of full geometries with all
their structure-preserving morphisms and the category of local models given
by primitive data and independently defined compatibility conditions
(Theorem~\ref{thm:8.18}). Objects are recovered up to isomorphism; morphisms
between fixed endpoints are recovered uniquely, with evaluation, composition,
and isomorphisms preserved (Corollary~\ref{cor:8.19}).

Here ``local'' refers to fragments of the data describing the geometry and
its maps, including the contexts and covers themselves. In Chapter~\ref{chap:2},
by contrast, the geometry and cover were fixed and the states were glued.
The reconstruction theorem applies to design descriptions divided among
teams when the required primitive data and compatibility conditions are
specified (\S\ref{sec:8.6}). In general the reconstruction uses compatible
data over all finite fragments. The special cases for lenses and protocols
explain when a finite table suffices (\S\ref{sec:8.7}); for finite models,
we also distinguish sufficient information from the checking procedure and
its computational cost (\S\ref{sec:8.8}).

\section*{Organization and reading guide}

The four parts comprise eight chapters. We assume basic familiarity with
sets, groups, rings, modules, categories, functors, and natural transformations.
The Preliminaries fix the notation; Chapter~\ref{chap:1} introduces AAT.
The chapters using sites, sheaves, cohomology, and idempotent completion
provide the needed definitions and references.

Begin with the objects and maps of Chapter~\ref{chap:1}, especially the core,
local geometry, and three-level projections (\S\S\ref{sec:1.4}--\ref{sec:1.7}).
For software applications, also read the lens and protocol semantics and
their encoding (\S\S\ref{sec:1.8}--\ref{sec:1.10}). Then choose a route:

\begin{itemize}
\item For local consistency and diagnosis, continue through
Chapters~\ref{chap:2}--\ref{chap:3}.
\item For transport and comparison of construction routes, concentrate on
Chapters~\ref{chap:4}--\ref{chap:6}, using the obstruction constructions of
Chapter~\ref{chap:2} and the diagnostic comparisons of Chapter~\ref{chap:3}
where they are invoked.
\item For classification and reconstruction, follow the construction of the
generated comparison in \S\ref{sec:5.10}, then read
Chapters~\ref{chap:6}--\ref{chap:8}, returning to the earlier transport and
base-change constructions as needed.
\end{itemize}

Related Work compares the objects, assumptions, admissible morphisms, and
conclusions with those of the closely related literature. Appendix~\ref{app:A}
lists the Lean declarations corresponding to the definitions, constructions,
and results, together with their assumptions and fixed source versions.
Appendix~\ref{app:B} gives the source locations and build instructions for
checking the proofs.
